\documentclass{MuTemp}
\usepackage[utf8]{inputenc}
\usepackage[T1]{fontenc}
\usepackage{geometry}
\usepackage{graphicx}
\usepackage{float}
\usepackage{hyperref}
\usepackage[english,indonesian]{babel}

\title{Optimasi Pemrosesan Template LaTeX Menggunakan Pendekatan Berbasis AI}
\author{Academic Support Agent}
\date{\today}

\begin{document}
\maketitle

\begin{abstract}

\end{abstract}

\chapter{PENDAHULUAN}


\chapter{TINJAUAN PUSTAKA}


% Original content (commented out):
% \documentclass[MuTemp]{MuTemp}

%%%%%%%%%%%%%%%%%%%%%%%%%%%%%%%%%%%%%%%%%%%%%%%%%%%%%%%
% Packages
\usepackage{array}
\usepackage{graphicx}
\usepackage[utf8]{inputenc}

%%%%%%%%%%%%%%%%%%%%%%%%%%%%%%%%%%%%%%%%%%%%%%%%%%%%%%%
% Header
\title{ABSRACT TITLE}
\author[*1]{Presenting Author}
\author[1,2]{Second Author}
\author[2]{Third Author}

\affil[1]{Department name, Institution 1} 
\affil[2]{Department name, Institution 2} 
% corresponding/presenting author email
\corrauthor{correspondingauthor@email.com}

%%%%%%%%%%%%%%%%%%%%%%%%%%%%%%%%%%%%%%%%%%%%%%%%%%%%%%%
\begin{document}

\maketitle

Enter your abstract here. Ensure that your document is no longer than 1 page in length. It is a good idea to start your abstract with some background and motivation, then describe your approach and summarize the key results.

A figure is not mandatory, but an eye-catching graphic is always good to include if it helps getting your message across. Just make sure that the total length of the abstract is \textbf{no more than 1 page} (including references).

\begin{figure}[htpb!]
\begin{center}
\includegraphics[scale=0.4]{sticker-1.png}
\caption{Sample figure.}
\label{Fig}
\end{center}
\end{figure}

Any references should appear together at the end of the abstract, listed alphabetically. Cite in the text using commands \citep{gargett1984,lilly1983} or \cite{multipleyeDMP2024}. 

%%%%%%%%%%%%%%%%%%%%%%%%%%%%%%%%%%%%%%%%%%%%%%%%%%%%%%%
\bibliographystyle{MuTemp}
\bibliography{MuTemp}

\end{document}

\end{document}